\section{问题重述}

光伏微网通过购电与储能共同满足负载。电池将低价或富余电量转移至后续时段，其容量、功率和效率决定可调节范围；采购不足与合同修订则产生额外费用。本文以实际结算账单为依据，协调采购承诺与储能运行。

四问逐步增加信息不确定性和决策机会：问题一求确定性单日联合调度；问题二预测负载与光伏后制定日前计划，缺口以五倍价格补购；问题三在规定时点利用新预报付费修订合同；问题四引入动态电价，分别推广日前与滚动采购。

图\ref{fig:modeling-framework}概括全文路线：以统一物理与结算模型为基础，构造分问采购策略，通过实际执行与时点消融确定最终配置。

\begin{figure}[H]\centering
% Native vector figure. Requires tikz and libraries arrows.meta,calc.
% Use inside a figure environment: \centering% Native vector figure. Requires tikz and libraries arrows.meta,calc.
% Use inside a figure environment: \centering% Native vector figure. Requires tikz and libraries arrows.meta,calc.
% Use inside a figure environment: \centering\input{figures/modeling_framework}
\begingroup
\definecolor{MMFInk}{HTML}{273B48}
\definecolor{MMFLine}{HTML}{566570}
\definecolor{MMFPale}{HTML}{F1F4F5}
\begin{tikzpicture}[
  x=1cm,y=1cm,
  font=\fontsize{9}{12}\selectfont,
  ink/.style={text=MMFInk},
  block/.style={ink,draw=MMFLine,line width=.55pt,fill=white,
    align=center,inner xsep=3pt,inner ysep=4pt,text width=3.23cm,
    minimum width=3.5cm},
  foundation/.style={block,fill=MMFPale,minimum height=1.5cm},
  question/.style={block,minimum height=2cm},
  outcome/.style={block,minimum height=1.5cm},
  flow/.style={draw=MMFLine,line width=.6pt,-{Stealth[length=1.8mm,width=1.3mm]}},
  bus/.style={draw=MMFLine,line width=.6pt}
]
  % Common units, information boundary, physical model and settlement.
  \node[foundation] (input) at (-5.625,0)
    {\textbf{附件输入}\\负荷、光伏功率与电价\\已发布的日内预报};
  \node[foundation] (units) at (-1.875,0)
    {\textbf{量纲与信息边界}\\10分钟功率转电量\\仅用决策时可得信息};
  \node[foundation,text width=6.98cm,minimum width=7.25cm] (mother) at (3.75,0)
    {\textbf{统一母模型}\\供给$\geq$净需求；储能容量、功率与效率约束\\计划量全额结算；调整与紧急补购单独计费};
  \draw[flow] (input.east) -- (units.west);
  \draw[flow] (units.east) -- (mother.west);

  % The four question settings share the mother model.
  \node[question] (q1) at (-5.625,-2.55)
    {\textbf{问题一：确定性调度}\\已知负荷、光伏与电价\\购电与储能联合LP\\确定性最优调度基准};
  \node[question] (q2) at (-1.875,-2.55)
    {\textbf{问题二：日前采购}\\整日历史残差构造场景\\情景LP制定共同计划\\计划后按实际供需执行};
  \node[question] (q3) at (1.875,-2.55)
    {\textbf{问题三：滚动采购}\\固定规则回放候选\\0、6、12、18时筛选\\更新剩余购电合同};
  \node[question] (q4) at (5.625,-2.55)
    {\textbf{问题四：动态电价}\\同日配对价格与净负荷\\日前版沿用母模型\\滚动版沿用候选筛选};
  \coordinate (split) at (3.75,-1.15);
  \draw[bus] (mother.south) -- (split);
  \draw[bus] (-5.625,-1.15) -- (5.625,-1.15);
  \foreach \n in {q1,q2,q3,q4}
    \draw[flow] (\n.north |- split) -- (\n.north);

  % Q1 is a joint deterministic LP, not a run of the frozen rule executor.
  % The closed-loop execution and screening evaluation apply to Q2--Q4.
  \coordinate (collect) at (0,-4.02);
  \foreach \n in {q2,q3,q4}
    \draw[bus] (\n.south) -- (\n.south |- collect);
  \draw[bus] (-5.625,-4.02) -- (5.625,-4.02);
  \node[ink,align=center,font=\fontsize{9}{12}\selectfont] at (1.3,-4.5)
    {\textbf{问题二至四：闭环执行与开发期选型}};

  \node[outcome] (execute) at (-5.625,-5.65)
    {\textbf{固定实时执行器}\\每10分钟依实际供需\\充放电与紧急补购};
  \node[outcome] (evaluate) at (-1.875,-5.65)
    {\textbf{连续状态与实际账单}\\储电量跨日连续推进\\计划、调整及紧急费};
  \node[outcome] (ablate) at (1.875,-5.65)
    {\textbf{开发期配置对照}\\滚动：四种筛选时点\\费用与紧急补购权衡};
  \node[outcome,fill=MMFPale] (final) at (5.625,-5.65)
    {\textbf{最终分问配置}\\日前两组保留母模型\\滚动两组全时点筛选};
  \draw[flow] (-5.625,-4.02) -- (execute.north);
  \draw[flow] (execute.east) -- (evaluate.west);
  \draw[flow] (evaluate.east) -- (ablate.west);
  \draw[flow] (ablate.east) -- (final.west);
\end{tikzpicture}
\endgroup

\begingroup
\definecolor{MMFInk}{HTML}{273B48}
\definecolor{MMFLine}{HTML}{566570}
\definecolor{MMFPale}{HTML}{F1F4F5}
\begin{tikzpicture}[
  x=1cm,y=1cm,
  font=\fontsize{9}{12}\selectfont,
  ink/.style={text=MMFInk},
  block/.style={ink,draw=MMFLine,line width=.55pt,fill=white,
    align=center,inner xsep=3pt,inner ysep=4pt,text width=3.23cm,
    minimum width=3.5cm},
  foundation/.style={block,fill=MMFPale,minimum height=1.5cm},
  question/.style={block,minimum height=2cm},
  outcome/.style={block,minimum height=1.5cm},
  flow/.style={draw=MMFLine,line width=.6pt,-{Stealth[length=1.8mm,width=1.3mm]}},
  bus/.style={draw=MMFLine,line width=.6pt}
]
  % Common units, information boundary, physical model and settlement.
  \node[foundation] (input) at (-5.625,0)
    {\textbf{附件输入}\\负荷、光伏功率与电价\\已发布的日内预报};
  \node[foundation] (units) at (-1.875,0)
    {\textbf{量纲与信息边界}\\10分钟功率转电量\\仅用决策时可得信息};
  \node[foundation,text width=6.98cm,minimum width=7.25cm] (mother) at (3.75,0)
    {\textbf{统一母模型}\\供给$\geq$净需求；储能容量、功率与效率约束\\计划量全额结算；调整与紧急补购单独计费};
  \draw[flow] (input.east) -- (units.west);
  \draw[flow] (units.east) -- (mother.west);

  % The four question settings share the mother model.
  \node[question] (q1) at (-5.625,-2.55)
    {\textbf{问题一：确定性调度}\\已知负荷、光伏与电价\\购电与储能联合LP\\确定性最优调度基准};
  \node[question] (q2) at (-1.875,-2.55)
    {\textbf{问题二：日前采购}\\整日历史残差构造场景\\情景LP制定共同计划\\计划后按实际供需执行};
  \node[question] (q3) at (1.875,-2.55)
    {\textbf{问题三：滚动采购}\\固定规则回放候选\\0、6、12、18时筛选\\更新剩余购电合同};
  \node[question] (q4) at (5.625,-2.55)
    {\textbf{问题四：动态电价}\\同日配对价格与净负荷\\日前版沿用母模型\\滚动版沿用候选筛选};
  \coordinate (split) at (3.75,-1.15);
  \draw[bus] (mother.south) -- (split);
  \draw[bus] (-5.625,-1.15) -- (5.625,-1.15);
  \foreach \n in {q1,q2,q3,q4}
    \draw[flow] (\n.north |- split) -- (\n.north);

  % Q1 is a joint deterministic LP, not a run of the frozen rule executor.
  % The closed-loop execution and screening evaluation apply to Q2--Q4.
  \coordinate (collect) at (0,-4.02);
  \foreach \n in {q2,q3,q4}
    \draw[bus] (\n.south) -- (\n.south |- collect);
  \draw[bus] (-5.625,-4.02) -- (5.625,-4.02);
  \node[ink,align=center,font=\fontsize{9}{12}\selectfont] at (1.3,-4.5)
    {\textbf{问题二至四：闭环执行与开发期选型}};

  \node[outcome] (execute) at (-5.625,-5.65)
    {\textbf{固定实时执行器}\\每10分钟依实际供需\\充放电与紧急补购};
  \node[outcome] (evaluate) at (-1.875,-5.65)
    {\textbf{连续状态与实际账单}\\储电量跨日连续推进\\计划、调整及紧急费};
  \node[outcome] (ablate) at (1.875,-5.65)
    {\textbf{开发期配置对照}\\滚动：四种筛选时点\\费用与紧急补购权衡};
  \node[outcome,fill=MMFPale] (final) at (5.625,-5.65)
    {\textbf{最终分问配置}\\日前两组保留母模型\\滚动两组全时点筛选};
  \draw[flow] (-5.625,-4.02) -- (execute.north);
  \draw[flow] (execute.east) -- (evaluate.west);
  \draw[flow] (evaluate.east) -- (ablate.west);
  \draw[flow] (ablate.east) -- (final.west);
\end{tikzpicture}
\endgroup

\begingroup
\definecolor{MMFInk}{HTML}{273B48}
\definecolor{MMFLine}{HTML}{566570}
\definecolor{MMFPale}{HTML}{F1F4F5}
\begin{tikzpicture}[
  x=1cm,y=1cm,
  font=\fontsize{9}{12}\selectfont,
  ink/.style={text=MMFInk},
  block/.style={ink,draw=MMFLine,line width=.55pt,fill=white,
    align=center,inner xsep=3pt,inner ysep=4pt,text width=3.23cm,
    minimum width=3.5cm},
  foundation/.style={block,fill=MMFPale,minimum height=1.5cm},
  question/.style={block,minimum height=2cm},
  outcome/.style={block,minimum height=1.5cm},
  flow/.style={draw=MMFLine,line width=.6pt,-{Stealth[length=1.8mm,width=1.3mm]}},
  bus/.style={draw=MMFLine,line width=.6pt}
]
  % Common units, information boundary, physical model and settlement.
  \node[foundation] (input) at (-5.625,0)
    {\textbf{附件输入}\\负荷、光伏功率与电价\\已发布的日内预报};
  \node[foundation] (units) at (-1.875,0)
    {\textbf{量纲与信息边界}\\10分钟功率转电量\\仅用决策时可得信息};
  \node[foundation,text width=6.98cm,minimum width=7.25cm] (mother) at (3.75,0)
    {\textbf{统一母模型}\\供给$\geq$净需求；储能容量、功率与效率约束\\计划量全额结算；调整与紧急补购单独计费};
  \draw[flow] (input.east) -- (units.west);
  \draw[flow] (units.east) -- (mother.west);

  % The four question settings share the mother model.
  \node[question] (q1) at (-5.625,-2.55)
    {\textbf{问题一：确定性调度}\\已知负荷、光伏与电价\\购电与储能联合LP\\确定性最优调度基准};
  \node[question] (q2) at (-1.875,-2.55)
    {\textbf{问题二：日前采购}\\整日历史残差构造场景\\情景LP制定共同计划\\计划后按实际供需执行};
  \node[question] (q3) at (1.875,-2.55)
    {\textbf{问题三：滚动采购}\\固定规则回放候选\\0、6、12、18时筛选\\更新剩余购电合同};
  \node[question] (q4) at (5.625,-2.55)
    {\textbf{问题四：动态电价}\\同日配对价格与净负荷\\日前版沿用母模型\\滚动版沿用候选筛选};
  \coordinate (split) at (3.75,-1.15);
  \draw[bus] (mother.south) -- (split);
  \draw[bus] (-5.625,-1.15) -- (5.625,-1.15);
  \foreach \n in {q1,q2,q3,q4}
    \draw[flow] (\n.north |- split) -- (\n.north);

  % Q1 is a joint deterministic LP, not a run of the frozen rule executor.
  % The closed-loop execution and screening evaluation apply to Q2--Q4.
  \coordinate (collect) at (0,-4.02);
  \foreach \n in {q2,q3,q4}
    \draw[bus] (\n.south) -- (\n.south |- collect);
  \draw[bus] (-5.625,-4.02) -- (5.625,-4.02);
  \node[ink,align=center,font=\fontsize{9}{12}\selectfont] at (1.3,-4.5)
    {\textbf{问题二至四：闭环执行与开发期选型}};

  \node[outcome] (execute) at (-5.625,-5.65)
    {\textbf{固定实时执行器}\\每10分钟依实际供需\\充放电与紧急补购};
  \node[outcome] (evaluate) at (-1.875,-5.65)
    {\textbf{连续状态与实际账单}\\储电量跨日连续推进\\计划、调整及紧急费};
  \node[outcome] (ablate) at (1.875,-5.65)
    {\textbf{开发期配置对照}\\滚动：四种筛选时点\\费用与紧急补购权衡};
  \node[outcome,fill=MMFPale] (final) at (5.625,-5.65)
    {\textbf{最终分问配置}\\日前两组保留母模型\\滚动两组全时点筛选};
  \draw[flow] (-5.625,-4.02) -- (execute.north);
  \draw[flow] (execute.east) -- (evaluate.west);
  \draw[flow] (evaluate.east) -- (ablate.west);
  \draw[flow] (ablate.east) -- (final.west);
\end{tikzpicture}
\endgroup

\caption{四问建模框架与最终策略形成过程}
\label{fig:modeling-framework}
\end{figure}

\section{模型假设}

\begin{enumerate}
 \item 假设附件功率为以标记时刻为右端点的10分钟平均功率，以$\Delta t=1/6$小时换算电量。
 \item 充电、放电效率分别为$\eta_c=\eta_d=0.9$，往返效率为0.81；忽略线路损耗和自放电。电池外侧充放电功率均不超过5000 kW，储电量保持在1200--10800 kWh。
 \item 不向外网售电，富余电量可免费弃置。原计划全额付费，取消部分不退款；另设退款解释作敏感性对照。
 \item 每10分钟观测当期负载、光伏、电价及储电量，紧急购电补足剩余缺口，外网补购量可按需获得；采购决策仅使用当时已观测或已发布的信息。
 \item 问题一首末储电量均为6000 kWh；其余问题自1月1日起逐日连续推进储电量。
\end{enumerate}

\section{符号说明与数据准备}

\subsection{主要符号与统一单位}

购电、充放电和储电量均以kWh计，电价以元/kWh计。以$p$表示电价、$c$表示充电量，避免成本与充电变量混用。跨日问题增加日期下标$D$，不致混淆时省略。

\begin{table}[htbp]\centering\small
\caption{主要符号及含义}\label{tab:symbols}
\begin{tabularx}{\linewidth}{p{0.19\linewidth}Xl}\toprule
符号 & 含义 & 单位\\\midrule
$D,t,s$ & 日期、时段$t=0,\ldots,143$、场景索引 & ---\\
$L_{D,t},R_{D,t}$ & 实际负载与光伏功率 & kW\\
$n_{D,t},p_{D,t}$ & 实际净负荷电量、实际交易电价 & kWh、元/kWh\\
$g_{D,t},a_{D,t}$ & 零时计划、最终生效常规购电量 & kWh\\
$c_t,d_t,e_t$ & 电池外侧充电量、放电量、时段初储电量 & kWh\\
$h_t,w_t$ & 紧急购电量、弃置电量 & kWh\\
$\pi_s,Q_s$ & 场景概率、场景剩余时域紧急费用 & ---、元\\
$\rho_t,\lambda$ & 参考备用量、CVaR权重 & kWh、---\\
$\beta_T,k_\rho$ & 终端偏离惩罚系数、备用缩放系数 & 元/kWh、---\\
$E_{\min},E_{\max},U$ & 储电量下上限、单段充放电量上限$5000/6$ & kWh\\
\bottomrule\end{tabularx}
\end{table}

\subsection{统一物理时间与模板列序}

每日保留144个10分钟时段，物理区间及净负荷定义为
\begin{equation}
 \begin{cases}
 I_t=[t/6,(t+1)/6),\\[2pt]
 n_{D,t}=(L_{D,t}-R_{D,t})\Delta t.
 \end{cases}
 \label{eq:time_energy}
\end{equation}
源标记00:10对应物理00:00--00:10，次日00:00对应当日23:50--24:00。模型按物理自然日结算；附件5保留原始首末标签，计算结果按源样本列序写入。

\begin{table}[htbp]\centering\small
\caption{计算时间与原模板显示的对应}\label{tab:physical_mapping}
\begin{tabularx}{\linewidth}{lLLL}\toprule
索引 & 源采样右端点 & 模型物理区间 & 原模板显示区间\\\midrule
$t=0$ & 00:10 & 00:00--00:10 & 0:10--0:20\\
$t=60$ & 10:10 & 10:00--10:10 & 10:10--10:20\\
$t=143$ & 次日00:00 & 23:50--24:00 & 次日0:00--0:10\\
\bottomrule\end{tabularx}
\end{table}

题目指定10:00--10:10取$t=60$，四小时区间累加连续24段，同日相邻正紧急量合并为一个区间。正文图表与工作簿共用这一物理索引，映射文件见附录。

\subsection{评价数据及图表口径}

问题一使用附件1。问题二至四的评价和填报覆盖2025年2月1日至12月31日334天，一月用于历史积累与有限候选校准。各问题的具体预测及可用信息分别在对应模型中说明。同一问题的比较策略从共同2月初态出发，各自连续执行；不同问题保留各自已冻结一月运行形成的初态。

预测器与实时执行器沿用一月冻结配置，候选筛选的适用问题及启用时点依据开发期回测确定。最终日前两组保留母模型，滚动两组在零时及全部合法日内修订时点启用筛选。

实际总费用由计划、调整及紧急三项账单组成，辅助终端罚项和风险权重用于优化评分。正文表格与附件5均从对应执行数组生成。热力图展示时序或策略差异，瀑布图分解费用变化，分位阴影带展示场景或历史窗口的经验分布。
